% Muestra tipográfica: reúne todos los elementos de la maqueta en pocas páginas
% para revisar el diseño. No forma parte del libro.
% Compilar desde esta carpeta con: latexmk muestra.tex
\documentclass[
  fontsize=10.5pt,
  twoside=true,
  headings=standardclasses,
  numbers=noenddot,
]{scrbook}

\usepackage{estilo/armaduras}
\addbibresource{bibliografia.bib}

\begin{document}

\setcounter{chapter}{4}
\chapter{La rigidez de una barra}

\begin{objetivos}
  \item Calcular la matriz de rigidez de una barra en coordenadas locales.
  \item Transformarla a coordenadas globales con la matriz $\mT$.
  \item Verificar que el resultado es simétrico y singular.
\end{objetivos}

Una barra de armadura solo trabaja a fuerza axial. Si conocemos su
\termino{módulo de elasticidad}{modulus of elasticity} $E$, el área $A$ de su
sección y su longitud $L$, la relación entre las fuerzas y los desplazamientos
de sus extremos, en la dirección de la barra, es
\nota{Esta es la ley de Hooke escrita para los dos extremos de la barra.}
\begin{equation}\label{eq:rigidez-local}
  \mkl = \EAL
  \begin{bmatrix} 1 & -1 \\ -1 & 1 \end{bmatrix}.
\end{equation}

\begin{recuerda}
La matriz de transformación lleva los desplazamientos globales a la dirección
de la barra: $\vul = \mT\,\vu$, con
$\mT = \begin{bmatrix} c & s & 0 & 0 \\ 0 & 0 & c & s \end{bmatrix}$.
Por lo tanto, $\mk = \mT\transp \mkl \mT$.
\end{recuerda}

\section{La barra del ejemplo A}

La \cref{fig:ejemplo-a} muestra la armadura del ejemplo A. Sus tres barras son
de acero ($E = \SI{200}{\giga\pascal}$) con un área de \SI{20}{\centi\metre\squared}.

\begin{figure}[ht]
  \centering
  \begin{tikzpicture}[scale=0.85]
    \coordinate (n1) at (0,0);
    \coordinate (n2) at (6,0);
    \coordinate (n3) at (3,4);
    \draw[barra] (n1) -- node[etiqueta barra] {1} (n2);
    \draw[barra] (n2) -- node[etiqueta barra] {2} (n3);
    \draw[barra] (n1) -- node[etiqueta barra] {3} (n3);
    \apoyoarticulado{(n1)}
    \apoyorodillo{(n2)}
    \cargapuntual{(n3)}{(1.4,0)}{\SI{20}{\kilo\newton}}
    \cargapuntual{(n3)}{(0,-1.4)}{\SI{50}{\kilo\newton}}
    \foreach \n/\pos in {n1/above left, n2/above right, n3/right}
      \node[nodo] at (\n) {};
    \node[etiqueta nodo, above left=2pt] at (n1) {1};
    \node[etiqueta nodo, above right=2pt] at (n2) {2};
    \node[etiqueta nodo, above right=4pt] at (n3) {3};
    \draw[eje] (-1.4,-0.9) -- ++(0.8,0) node[right, font=\small] {$x$};
    \draw[eje] (-1.4,-0.9) -- ++(0,0.8) node[above, font=\small] {$y$};
  \end{tikzpicture}
  \caption{Ejemplo A: armadura triangular isostática.}
  \label{fig:ejemplo-a}
\end{figure}

\begin{python}
En Python, el producto de matrices se escribe con el operador \py{@}.
La transpuesta de un arreglo \py{T} de NumPy es \py{T.T}.
\end{python}

La función que construye la matriz de una barra traduce la
\cref{eq:rigidez-local} casi literalmente:

\begin{codigo}
import numpy as np

def rigidez_barra(E, A, L, c, s):
    """Matriz de rigidez global (4 x 4) de una barra."""
    k_local = E * A / L * np.array([[1, -1], [-1, 1]])
    T = np.array([[c, s, 0, 0], [0, 0, c, s]])
    return T.T @ k_local @ T
\end{codigo}

\begin{cuidado}
La matriz de una barra aislada es singular: la barra puede moverse como cuerpo
rígido. Intentar resolver un sistema con ella produce un error. Los apoyos son
los que eliminan esa singularidad.
\end{cuidado}

\begin{table}[ht]
  \centering
  \caption{Geometría de las barras del ejemplo A.}
  \begin{tabular}{@{}c S[table-format=1.2] S[table-format=-1.2] S[table-format=-1.2]@{}}
    \toprule
    {Barra} & {$L$ (\si{\metre})} & {$c$} & {$s$} \\
    \midrule
    1 & 6.00 & 1.00 & 0.00 \\
    2 & 5.00 & -0.60 & 0.80 \\
    3 & 5.00 & 0.60 & 0.80 \\
    \bottomrule
  \end{tabular}
\end{table}

\begin{parasabermas}
La deducción completa de la matriz de rigidez se encuentra en
\textcite{kassimali2012} y en \textcite{mcguire2000}.
\end{parasabermas}

\section*{Ejercicios}

\begin{ejercicio}
Calcular la matriz de rigidez en coordenadas globales de la barra~2 del
ejemplo A y verificar que es simétrica.
\end{ejercicio}

\puntodecontrol{05}

\printbibliography[heading=subbibliography]

\end{document}
